\section{Discussion}
\label{sec:discussion}

\paragraph{Two complementary failure modes}
On short sharded tasks, retaining the model's own replies lowered final
success, replicating the direction of \citet{huang2026benefit} under a frozen
paired analysis. Over 100 turns of state tracking the loss was far larger:
accuracy fell to about a sixth of its initial level within twenty turns, in
every scenario and for three models of different families and sizes, and a
compact verified state removed most of it. We do not claim that the short-task
effect grows causally into the long-horizon collapse, and Study~2 cannot
separate history length, the number of state updates and error accumulation,
because the verified summary changes all three. The practical finding stands:
in stateful multi-turn use, keeping the raw transcript as the only memory has a
large and avoidable accuracy cost.

\paragraph{Why quality metrics stay silent}
Consistency, coherence and stability describe local properties of responses. A
model that has adopted a wrong state can reason coherently and repeatably from
it. In our data these metrics stayed level or rose as accuracy fell, suggesting
that accumulated context makes responses more anchored rather than more
erratic, in line with reports of self-conditioning and sycophancy
\citep{sinha2025illusion,sharma2024sycophancy}. The failure mode does not
produce the incoherence or disagreement these metrics register. Two classes of
signal remain open. The first is uncertainty. Token entropy, the signal used
by ERGO \citep{khalid2025ergo}, needs token probabilities, which our protocols
did not collect. Semantic entropy \citep{farquhar2024detecting} does not: its
discrete form clusters sampled answers by meaning and needs only the samples.
Our consistency metric is related, because it also summarises agreement among
the $K$ sampled answers, but it is a pairwise agreement rate, not an entropy,
and we did not compute semantic entropy. Both semantic and token-level entropy
therefore remain untested here. The second class is counterfactual signals,
which compare the deployed context with a history-free or state-grounded
re-query and so target the anchoring directly. CRL tracked the
current gap rather than anticipating it, and comparing answers rather than
readiness is a natural next step.

\paragraph{Implications for LLM-based information systems}
Four design rules follow. \emph{Keep state outside the transcript}: where the
state can be represented explicitly, as in ledgers, inventories, protocols and
schedules, maintain it in a verified store and send the current state with each
request. In our setting this kept accuracy 0.52--0.59 above full history and
shrank the prompt at turn 100 from tens of thousands of tokens to a few
hundred. Our solver-verified summary is an upper reference; model-written
summaries should be expected to do less well. \emph{Do not rely on
single-stream quality scores}: a dashboard of one conversation's consistency,
coherence or stability will typically remain reassuring while accuracy
collapses. \emph{Monitor by comparison or aggregation}: pooling concurrent
sessions makes realistic drops detectable, and counterfactual re-queries target
the history effect itself, at the price of delayed per-session alarms or extra
calls. \emph{Test the monitor before trusting it}: verify that each detector
can fire on streams of the intended length and that its reference period is
free of the degradation to be detected.

\paragraph{Limitations}
Both studies use arithmetic and state-tracking tasks with exact answers, which
make accuracy observable but are narrower than open-ended dialogue. All models
were accessed through one academic gateway, the two studies use different
model sets, and newer models may degrade differently. The models are not the
strongest available. This is partly a property of the design: per-model
calibration needs an informative accuracy range, and a stronger open model,
Gemma~4 31B, stayed above the hard band even at the contingency levels of the
difficulty ladder and could not be calibrated (Supplementary
Section~S2.2). Testing frontier models therefore
requires a harder ladder, not only more calls. Study~1 estimates a net
effect that includes length and position differences, and Study~2's control is
a packaged intervention, not a mechanism test. Accuracy, robustness and
efficiency were unobservable before the final turn in Study~1. Study~2's CS, RS
and SS samples are complete only for GPT-4o-mini. Monitoring analyses beyond
the frozen tests are exploratory or post-hoc, two deviations for Llama-3 70B
followed gate results and are disclosed, and one investigator performed the
task audit and calibration.
